\subsection{Anvendelse af Boltzmann ligning}\label{Braginskii example}
Givet 2D rummet, med fordelingensfunktionen:

\begin{align*}
    &\beta = (1, 2)\\
    &\mathbf{x} = (x_1, x_2)\\
    &\mathbf{v} = (v_1, v_2)\\
    &f_a(t, x_1, x_2, v_1, v_2) = t x_1^2 + x_2 v_1 + v_2^2
\end{align*}

Her ses bort fra kraftleddet ved at sætte $\frac{F_{a\beta}}{m_a} = \alpha$, 

Ligningen er er givet ved \ref{eq:boltzmann_plasma}. For Einstein summation over $\beta$ gælder:
\begin{equation*}
    \partial_{x_\beta}(v_\beta f_a) = \partial_{x_1}(v_1 f_a) + \partial_{x_2}(v_2 f_a).
\end{equation*}

Herved kan de afledte findes ved:
\begin{itemize}
    \item Tidsafledte
\end{itemize}
\begin{equation*}
    \partial_t f_a = \partial_t(t x_1^2 + x_2 v_1 + v_2^2) = x_1^2.
\end{equation*}

\begin{itemize}
    \item Rumleddet
\end{itemize}

\begin{equation*}
    \partial_{x_\beta}(v_\beta f_a) = \partial_{x_1}(v_1 f_a) + \partial_{x_2}(v_2 f_a) = 2v_1tx_1 + v_1v_2
\end{equation*}

\begin{itemize}
    \item Hastighedsleddet
\end{itemize}
\begin{equation*}
    \partial_{v_\beta}(\alpha f_a) = \partial_{v_1}(\alpha f_a) + \partial_{v_2}(\alpha f_a) = \alpha x_2 + 2\alpha v_2
\end{equation*}
\begin{itemize}
    \item Samling med kollisionsleddet
\end{itemize}
\begin{equation*}
    x_1^2 + 2t x_1 v_1 + v_1 v_2 + \alpha x_2 + 2\alpha v_2 = C_a
\end{equation*}
